% SPDX-License-Identifier: Apache-2.0
\section{The Scanout Beside the Video Peripheral}
\label{sec:x-scanvideo}

\S\ref{sec:x-scanpair} has the scanout's two ends. This section puts
the pixel end beside the video peripheral, the one the board's third
slot already holds, and lets a host program it over that slot. The
result, \code{ScanVideo}, is one lowered unit of seven:

\begin{center}\footnotesize
\begin{tabular}{@{}l@{}}
\code{slot -- LiteSplit -- Hdmi, below 0x80 -- VidMux -- pins} \\
\code{slot -- LiteSplit -- ScanCtl, from 0x80 -- LinePair -- VidMux} \\
\code{Raster -- LinePair -- ScanTap -- ScanCtl} \\
\end{tabular}
\end{center}

\code{LiteSplit} divides the slot by address bit 7: \code{Hdmi}'s
words stay below \code{0x80}, as they were, and \code{ScanCtl}'s
registers start there. \code{ScanCtl} holds the frame's base, the bit
that shows the scanout rather than the framebuffer, and the underflow
bit, which a write clears. \code{VidMux} puts one picture or the other
on the pins, with \code{Hdmi}'s syncs either way.

\subsection{Why the beam is counted twice}

A wire between two units is a wire in the netlist, read in the same
cycle it is driven. A run has to step one unit and then the other,
so a wire is read right only by a unit that runs after the one
driving it. Here the control goes one way, from the registers to the
line pair, and the status comes back, so whichever order the units
run in, one direction is read a step late. The run and the netlist
would then disagree, and the co-simulation says so.

Two things break the loop. The underflow bit comes back over a
channel, through \code{ScanTap}: a channel commits at the end of the
step, whichever unit ran first, so the bit reaches the registers a
few cycles late in the run and the netlist alike. And the line pair
does not read the beam out of \code{Hdmi}: \code{Raster} counts it
again, with \code{Hdmi}'s own functions and parameters, so it runs
first and every wire it drives is read after it.

Two counters that must agree are a risk, so the run checks them on
every cycle, through a reset in the middle of a frame:

\lstinputlisting[rangebeginprefix=//\ begin\{,rangebeginsuffix=\},%
rangeendprefix=//\ end\{,rangeendsuffix=\},includerangemarker=false,%
linerange=raster-raster,caption={\code{lib/parts/src/hdmi.rs}: the beam,
counted again.},label={lst:x-scanvideo-raster}]{lib/parts/src/hdmi.rs}

\subsection{What the run checks}

The host paints one pixel of the framebuffer, through \code{Hdmi}'s
cursor below \code{0x80}, and frame 0 shows it and black elsewhere.
Then it writes the frame's base and sets the scanout bit, and frames
1 and 2 show every visible pixel as the word at its place in memory.
Frame 3's memory stalls for a whole line, and a read of the status
after it finds the underflow bit set; a write clears it at the start
of frame 4, and a read at its end finds it clear.

\begin{figure}[htbp]
\centering
\input{scanvideo_timing.tex}
\caption{Frame 3: no word arrives while the memory holds the row
after, \code{rgot} is short at that row's start, \code{under} sets in
the pair, and \code{seen} follows it in the registers a few cycles
later, over the channel.}
\label{fig:x-scanvideo}
\end{figure}

The entire \code{ScanVideo} module compiles to a single netlist,
simulated against the run under nvc and Verilator.

\lstinputlisting[rangebeginprefix=//\ begin\{,rangebeginsuffix=\},%
rangeendprefix=//\ end\{,rangeendsuffix=\},includerangemarker=false,%
linerange=video_run-video_run,caption={\code{lib/parts/src/scanout.rs}:
the seven units, in the order that puts every driver before its
reader.},label={lst:x-scanvideo-run}]{lib/parts/src/scanout.rs}

\lstinputlisting[caption={\code{ex\_scanvideo.rs}. Run with \code{bazel run //lib/examples:ex\_scanvideo}.},label={lst:x-scanvideo}]{lib/examples/ex_scanvideo.rs}

\lstinputlisting[style=ascii,caption={What \code{ex\_scanvideo} printed when the build ran it: the pixels checked, the underflow bit read over the slot, the cycles the two counts agreed, and the start of the netlist.},label={lst:x-scanvideo-out}]{out_scanvideo.txt}
